\ifthenelse{\equal{\lang}{ko}}{%
본 장에서는 법률·정부 정보 접근을 위한 기존 접근 방식을 검토하고, 그 한계를 강조하며, 다국어 NLP/LLM의 발전을 정리하고, \textbf{포용적 접근(inclusive access)}을 위한 다국어 AI 프레임워크의 필요성을 명확히 한다.

\subsection{공공 법률·정부 정보 접근 방식}
공공 법률 및 정부 정보에 대한 접근은 오랫동안 정보 격차의 원천이었다. \textbf{키워드 기반 검색과 정적 포털} (예: 정부 웹사이트, RSS)은 배포가 쉽지만 표면적 일치에 의존하여 불필요한 잡음을 발생시키거나 미묘한 의미를 놓치기 쉽다 \cite{zhong2020legalai}. \textbf{수작업 번역 및 상담}은 정확한 맥락을 제공하지만 느리고 비용이 많이 들며 확장성이 부족하다 \cite{aletras2016echr}. \textbf{정보 검색 데이터베이스} (예: 법령·판례 검색 시스템)는 방대한 자료를 통합하지만 다국어 지원이 제한적이고, 비전문가가 이해하기 어려운 형식으로 결과를 제공한다 \cite{niklaus2023multilegal}.

\begin{table}[h]
\centering
\small
\begin{tabular}{p{0.25\linewidth}p{0.65\linewidth}}
\toprule
접근 방식 & 한계 \\
\midrule
키워드 검색 & 맥락 부족, 높은 오탐률 \\
수작업 번역·상담 & 정확하나 비용 높고 느림 \\
정보 검색 DB & 제한된 다국어 지원, 비전문가 이해 어려움 \\
\bottomrule
\end{tabular}
\caption{기존 공공 정보 접근 방식의 한계.}
\end{table}

\subsection{법률·정책 텍스트를 위한 다국어 NLP와 LLM}
초기 연구들은 규칙 기반 조항 추출, 계약 분석, 법률 검색에 NLP를 적용하였다 \cite{chalkidis2020legalbert}. Transformer 기반 다국어 인코더(mBERT, XLM-R, mT5)의 등장은 언어 간 전이(cross-lingual transfer)를 통해 저자원 언어에서도 강건한 표현 학습을 가능케 하였다 \cite{devlin2019bert, conneau2020xlmr, xue2021mt5}. 최근 GPT-4, LLaMA-2와 같은 대규모 언어 모델(LLMs)은 강력한 의미 추론과 요약, 그리고 다국어 이해 능력을 보여준다 \cite{openai2023gpt4, touvron2023llama2}. 그러나 이러한 능력을 \emph{시민 대면(citizen-facing) 정보 접근}에 통합한 프레임워크는 여전히 드물다. LexGLUE, LegalBench, CUAD와 같은 벤치마크는 법률 NLP 과제를 표준화하였으나 \cite{chalkidis2022lexglue, guha2023legalbench, hendrycks2021cuad}, 대부분 영어 중심이며 다국어 포용성을 충분히 다루지 않는다.

\subsection{근거 기반 생성과 책임 있는 AI}
병행 연구에서는 사실성(factuality)과 투명성을 강조한다. 검색 증강 생성(RAG)은 외부 근거에 기반하여 생성의 신뢰성을 높이며 \cite{lewis2020rag}, 요약의 충실성(faithfulness)과 환각(hallucination) 탐지에 관한 연구는 LLM 출력의 사실 검증 필요성을 부각한다 \cite{maynez2020faithfulness, ji2023hallucination}. 동시에 Model Cards, 언어 모델의 사회적 위험 분류 등은 배포 시 해석 가능성과 윤리의 중요성을 강조한다 \cite{mitchell2021modelcards, bender2021stochastic, weidinger2021taxonomy}. 그러나 대부분은 시민 대면 접근성 시스템에 내재화되지 않았다.

\subsection{연구 공백}
현행 방법들은 관할권과 언어를 넘나드는 확장 가능하고, 최신성이 유지되며, 근거에 기반한 시민 대면 정보 접근을 제공하지 못한다. 더 나아가 사실성·투명성·인간 감독과 같은 안전장치는 시스템 설계에 거의 통합되지 않았다 \cite{gama2014drift}. 이는 본 연구의 동기를 제공한다: 법률 및 정부 지식에 대한 \textbf{포용적 접근}을 지원하는 다국어·근거 기반 LLM 프레임워크.

\begin{table}[t]
\centering
\small
\resizebox{\columnwidth}{!}{%
\begin{tabular}{p{2.2cm} p{3.3cm} p{5.0cm}}
\toprule
 & \textbf{법률 정보 검색} & \textbf{다국어 포용적 접근 (본 연구)} \\
\midrule
주요 초점 & 검색, 조항/요소 추출, 법률 QA & \textbf{변화 탐지} $\rightarrow$ \textbf{다국어 분류} $\rightarrow$ \textbf{근거 기반 요약·LLM 보고} (엔드투엔드) \\
시간성 & 정적/주기적 업데이트 & \textbf{준실시간} 변화 감지 및 드리프트 대응 \\
언어 & 주로 단일 언어(영어 중심) & \textbf{다국어} (영·한·불 + 저자원 스트레스 테스트) \\
대상 사용자 & 법률 전문가 & \textbf{비전문가·다국어 공동체} \\
핵심 방법 & 규칙/IR + 전통적 NLP & \textbf{LLM 추론}, 근거 기반 설명 가능성 \\
출력 & 문서/조각 & \textbf{출처 포함 구조화·근거 연결 요약} \\
\bottomrule
\end{tabular}
}
\caption{법률 정보 검색과 \emph{다국어 포용적 접근}의 개념적 차이.}
\label{tab:conceptual}
\end{table}


}{%
This section reviews existing approaches to accessing legal and government information, highlights their limitations, situates advances in multilingual NLP/LLMs, and clarifies the need for a multilingual AI framework for \textbf{inclusive access}.

\subsection{Approaches to Public Legal and Government Information}
Access to public legal and government information has long been a source of information inequality. \textbf{Keyword-based search and static portals} (e.g., government websites, RSS) are easy to deploy but rely on surface-level matches, often generating irrelevant noise or missing nuanced meaning \cite{zhong2020legalai}. \textbf{Manual translation and advisory services} provide accurate context but are slow, costly, and lack scalability \cite{aletras2016echr}. \textbf{Information-retrieval databases} (e.g., statute and case-law search systems) consolidate large corpora but offer limited multilingual support and return results in formats that non-experts find difficult to understand \cite{niklaus2023multilegal}.

\begin{table}[h]
\centering
\small
\begin{tabular}{p{0.25\linewidth}p{0.65\linewidth}}
\toprule
Approach & Limitation \\
\midrule
Keyword search & Misses context, high false positives \\
Manual translation/advice & Accurate but costly and slow \\
IR databases & Limited multilingual support, hard for non-experts \\
\bottomrule
\end{tabular}
\caption{Limitations of existing public-information access approaches.}
\end{table}

\subsection{Multilingual NLP and LLMs for Legal and Policy Texts}
Early work applied NLP to rule-based clause extraction, contract analysis, and legal retrieval \cite{chalkidis2020legalbert}. The emergence of transformer-based multilingual encoders (mBERT, XLM-R, mT5) enabled robust representation learning even for low-resource languages via cross-lingual transfer \cite{devlin2019bert, conneau2020xlmr, xue2021mt5}. Recent large language models (LLMs) such as GPT-4 and LLaMA-2 demonstrate strong semantic reasoning, summarization, and multilingual understanding \cite{openai2023gpt4, touvron2023llama2}. However, frameworks that integrate these capabilities into \emph{citizen-facing} information access remain scarce. Benchmarks such as LexGLUE, LegalBench, and CUAD have standardized legal NLP tasks \cite{chalkidis2022lexglue, guha2023legalbench, hendrycks2021cuad}, but most are English-centric and do not sufficiently address multilingual inclusiveness.

\subsection{Evidence-Grounded Generation and Responsible AI}
Parallel research emphasizes factuality and transparency. Retrieval-augmented generation (RAG) grounds generation in external evidence to improve reliability \cite{lewis2020rag}, while work on summarization faithfulness and hallucination detection highlights the need to fact-check LLM outputs \cite{maynez2020faithfulness, ji2023hallucination}. At the same time, Model Cards and taxonomies of social risks from language models stress the importance of interpretability and ethics in deployment \cite{mitchell2021modelcards, bender2021stochastic, weidinger2021taxonomy}. Yet most remain external to citizen-facing accessibility systems.

\subsection{Research Gap}
Current methods fail to deliver scalable, current, and evidence-grounded citizen-facing information access across jurisdictions and languages. Moreover, safeguards such as factuality, transparency, and human oversight are rarely integrated into system design \cite{gama2014drift}. This motivates our work: a multilingual, evidence-grounded LLM framework that supports \textbf{inclusive access} to legal and governmental knowledge.

\begin{table}[t]
\centering
\small
\resizebox{\columnwidth}{!}{%
\begin{tabular}{p{2.2cm} p{3.3cm} p{5.0cm}}
\toprule
 & \textbf{Legal Information Retrieval} & \textbf{Multilingual Inclusive Access (This Work)} \\
\midrule
Primary Focus & Retrieval, clause/element extraction, legal QA & \textbf{Change detection} $\rightarrow$ \textbf{multilingual classification} $\rightarrow$ \textbf{evidence-grounded summarization / LLM reporting} (end-to-end) \\
Temporality & Static/periodic updates & \textbf{Near real-time} change detection and drift-aware adaptation \\
Languages & Mostly monolingual (English-centric) & \textbf{Multilingual} (EN/KO/FR + low-resource stress tests) \\
Target Users & Legal experts & \textbf{Non-experts and multilingual communities} \\
Method Core & Rules/IR + classical NLP & \textbf{LLM reasoning}, evidence-linked explainability \\
Outputs & Documents/fragments & \textbf{Structured, evidence-linked summaries} with provenance \\
\bottomrule
\end{tabular}
}
\caption{Conceptual distinction between Legal Information Retrieval and \emph{Multilingual Inclusive Access}.}
\label{tab:conceptual}
\end{table}


}
